% Options for packages loaded elsewhere
\PassOptionsToPackage{unicode}{hyperref}
\PassOptionsToPackage{hyphens}{url}
\PassOptionsToPackage{dvipsnames,svgnames,x11names}{xcolor}
%
\documentclass[
  twocolumn,
  10pt]{article}

\usepackage{amsmath,amssymb}
\usepackage{iftex}
\ifPDFTeX
  \usepackage[T1]{fontenc}
  \usepackage[utf8]{inputenc}
  \usepackage{textcomp} % provide euro and other symbols
\else % if luatex or xetex
  \usepackage{unicode-math}
  \defaultfontfeatures{Scale=MatchLowercase}
  \defaultfontfeatures[\rmfamily]{Ligatures=TeX,Scale=1}
\fi
\usepackage{lmodern}
\ifPDFTeX\else  
    % xetex/luatex font selection
\fi
% Use upquote if available, for straight quotes in verbatim environments
\IfFileExists{upquote.sty}{\usepackage{upquote}}{}
\IfFileExists{microtype.sty}{% use microtype if available
  \usepackage[]{microtype}
  \UseMicrotypeSet[protrusion]{basicmath} % disable protrusion for tt fonts
}{}
\makeatletter
\@ifundefined{KOMAClassName}{% if non-KOMA class
  \IfFileExists{parskip.sty}{%
    \usepackage{parskip}
  }{% else
    \setlength{\parindent}{0pt}
    \setlength{\parskip}{6pt plus 2pt minus 1pt}}
}{% if KOMA class
  \KOMAoptions{parskip=half}}
\makeatother
\usepackage{xcolor}
\usepackage[margin=2cm]{geometry}
\setlength{\emergencystretch}{3em} % prevent overfull lines
\setcounter{secnumdepth}{-\maxdimen} % remove section numbering
% Make \paragraph and \subparagraph free-standing
\makeatletter
\ifx\paragraph\undefined\else
  \let\oldparagraph\paragraph
  \renewcommand{\paragraph}{
    \@ifstar
      \xxxParagraphStar
      \xxxParagraphNoStar
  }
  \newcommand{\xxxParagraphStar}[1]{\oldparagraph*{#1}\mbox{}}
  \newcommand{\xxxParagraphNoStar}[1]{\oldparagraph{#1}\mbox{}}
\fi
\ifx\subparagraph\undefined\else
  \let\oldsubparagraph\subparagraph
  \renewcommand{\subparagraph}{
    \@ifstar
      \xxxSubParagraphStar
      \xxxSubParagraphNoStar
  }
  \newcommand{\xxxSubParagraphStar}[1]{\oldsubparagraph*{#1}\mbox{}}
  \newcommand{\xxxSubParagraphNoStar}[1]{\oldsubparagraph{#1}\mbox{}}
\fi
\makeatother


\providecommand{\tightlist}{%
  \setlength{\itemsep}{0pt}\setlength{\parskip}{0pt}}\usepackage{longtable,booktabs,array}
\usepackage{calc} % for calculating minipage widths
% Correct order of tables after \paragraph or \subparagraph
\usepackage{etoolbox}
\makeatletter
\patchcmd\longtable{\par}{\if@noskipsec\mbox{}\fi\par}{}{}
\makeatother
% Allow footnotes in longtable head/foot
\IfFileExists{footnotehyper.sty}{\usepackage{footnotehyper}}{\usepackage{footnote}}
\makesavenoteenv{longtable}
\usepackage{graphicx}
\makeatletter
\newsavebox\pandoc@box
\newcommand*\pandocbounded[1]{% scales image to fit in text height/width
  \sbox\pandoc@box{#1}%
  \Gscale@div\@tempa{\textheight}{\dimexpr\ht\pandoc@box+\dp\pandoc@box\relax}%
  \Gscale@div\@tempb{\linewidth}{\wd\pandoc@box}%
  \ifdim\@tempb\p@<\@tempa\p@\let\@tempa\@tempb\fi% select the smaller of both
  \ifdim\@tempa\p@<\p@\scalebox{\@tempa}{\usebox\pandoc@box}%
  \else\usebox{\pandoc@box}%
  \fi%
}
% Set default figure placement to htbp
\def\fps@figure{htbp}
\makeatother

\usepackage{fontspec}
\setmainfont{Times New Roman}
\usepackage{fancyhdr}
\usepackage{booktabs}
\usepackage{graphicx}
\usepackage{caption}
\captionsetup{font=small,labelfont=bf}
\usepackage[none]{hyphenat}
\usepackage{titlesec}
\titleformat{\section}{\normalfont\fontsize{11}{13}\bfseries}{\thesection.}{0.5em}{\MakeUppercase}
\titlespacing*{\section}{0pt}{10pt}{4pt}
\pagestyle{fancy}
\fancyhf{}
\renewcommand{\headrulewidth}{0pt}
\fancyfoot[C]{\footnotesize Maestría en Inteligencia Artificial Aplicada, UIDE — Visión por Computador — Práctica Semana 2}
\setlength{\columnsep}{0.6cm}
\makeatletter
\@ifpackageloaded{caption}{}{\usepackage{caption}}
\AtBeginDocument{%
\ifdefined\contentsname
  \renewcommand*\contentsname{Table of contents}
\else
  \newcommand\contentsname{Table of contents}
\fi
\ifdefined\listfigurename
  \renewcommand*\listfigurename{List of Figures}
\else
  \newcommand\listfigurename{List of Figures}
\fi
\ifdefined\listtablename
  \renewcommand*\listtablename{List of Tables}
\else
  \newcommand\listtablename{List of Tables}
\fi
\ifdefined\figurename
  \renewcommand*\figurename{Figure}
\else
  \newcommand\figurename{Figure}
\fi
\ifdefined\tablename
  \renewcommand*\tablename{Table}
\else
  \newcommand\tablename{Table}
\fi
}
\@ifpackageloaded{float}{}{\usepackage{float}}
\floatstyle{ruled}
\@ifundefined{c@chapter}{\newfloat{codelisting}{h}{lop}}{\newfloat{codelisting}{h}{lop}[chapter]}
\floatname{codelisting}{Listing}
\newcommand*\listoflistings{\listof{codelisting}{List of Listings}}
\makeatother
\makeatletter
\makeatother
\makeatletter
\@ifpackageloaded{caption}{}{\usepackage{caption}}
\@ifpackageloaded{subcaption}{}{\usepackage{subcaption}}
\makeatother

\usepackage{bookmark}

\IfFileExists{xurl.sty}{\usepackage{xurl}}{} % add URL line breaks if available
\urlstyle{same} % disable monospaced font for URLs
\hypersetup{
  colorlinks=true,
  linkcolor={blue},
  filecolor={Maroon},
  citecolor={Blue},
  urlcolor={Blue},
  pdfcreator={LaTeX via pandoc}}


\author{}
\date{}

\begin{document}


\twocolumn[
\begin{center}
{\fontspec{Arial}\bfseries\fontsize{16}{19}\selectfont Detección automática de incendios: comparación de descriptores clásicos de textura y arquitecturas de aprendizaje profundo sobre imágenes reales}\\[8pt]
{\fontspec{Times New Roman}\fontsize{12}{14}\selectfont [Nombre Apellido 1], [Nombre Apellido 2], [Nombre Apellido 3] — Grupo 2}\textbackslash{[}4pt{]}
\{\fontspec{Times New Roman}\itshape\fontsize{10}{12}\selectfont Maestría en Inteligencia Artificial Aplicada, Universidad Internacional del Ecuador (UIDE), Quito, Ecuador\}
\textbackslash end\{center\} \vspace{10pt}

\noindent\textbf{\textit{Resumen}} --- Este trabajo compara descriptores
clásicos de textura (HOG, Gabor, GLCM/Haralick, LBP, SIFT/ORB) contra
arquitecturas de aprendizaje profundo (CNN entrenada desde cero,
transfer learning y fine-tuning con MobileNetV2) para la clasificación
de imágenes de incendios forestales, usando el conjunto de datos real
``Fire Dataset'' (Kaggle, 799 imágenes de entrenamiento, 200 de prueba).
Los resultados muestran que el transfer learning con MobileNetV2 alcanza
el mejor desempeño (91.00\% de accuracy), superando a la CNN entrenada
desde cero (84.00\%) y a los descriptores clásicos combinados con SVM
(84.00\%). Se documenta además un hallazgo relevante en el análisis de
textura: las regiones de fuego presentan mayor energía de Haralick
(mayor uniformidad de intensidad) que el fondo sin incendio, contrario a
la intuición inicial. Finalmente, se discute por qué el accuracy global
puede ser una métrica insuficiente en aplicaciones de seguridad, dado
que el costo de un falso negativo (incendio no detectado) supera
ampliamente al de un falso positivo.

\vspace{6pt}

\noindent\textbf{\textit{Palabras Clave}} --- visión por computador,
detección de incendios, HOG, GLCM, redes neuronales convolucionales,
transfer learning, MobileNetV2. \vspace{10pt} {]}

\subsection{Introducción y
metodología}\label{introducciuxf3n-y-metodologuxeda}

La detección temprana de incendios mediante visión por computador es un
caso de uso crítico donde el costo de los distintos tipos de error no es
simétrico: no detectar un incendio real (falso negativo) es
considerablemente más grave que activar una alerta innecesaria (falso
positivo). Este trabajo evalúa dos familias de enfoques ---descriptores
clásicos de textura/forma y aprendizaje profundo--- sobre el conjunto de
datos ``Fire Dataset'' de Kaggle (799 imágenes de entrenamiento, 200 de
prueba, 224×224 px, escala de grises), siguiendo el pipeline de la guía
de práctica de la Semana 2 (Sesión 1: reconocimiento de patrones; Sesión
2: redes neuronales convolucionales).

\subsection{Descriptores clásicos de textura y forma (Sesión
1)}\label{descriptores-cluxe1sicos-de-textura-y-forma-sesiuxf3n-1}

\textbf{HOG.} Al aumentar el tamaño de celda de 8×8 a 16×16 píxeles, el
vector descriptor se redujo de 26 244 a 6 084 elementos, consecuencia
directa de que el número de celdas de la malla se reduce a la cuarta
parte.

\textbf{Filtros de Gabor.} Un banco ampliado de 4 a 5 orientaciones
reveló una nueva activación direccional, confirmando la sensibilidad
angular de estos filtros.

\textbf{HOG + SVM.} El kernel RBF (84.00\%) superó marginalmente al
lineal (83.50\%), sugiriendo que el descriptor HOG ya es razonablemente
separable de forma lineal.

\textbf{PCA + K-means.} Al incrementar \emph{k} de 2 a 3, el algoritmo
fragmentó geométricamente el grupo mayoritario en dos subgrupos, sin
corresponder a una clase real ---evidencia de que K-means agrupa por
densidad espacial, no por semántica.

\textbf{GLCM/Haralick y LBP.} Se comparó la propiedad \emph{energy}
entre tres muestras reales del dataset (Tabla 1, Figura 1).

\begin{center}
\small
\begin{tabular}{@{}lc@{}}
\toprule
\textbf{Muestra} & \textbf{Energía (GLCM)} \\
\midrule
Fuego — muestra 1 & 0.0355 \\
Fuego — muestra 2 & 0.0274 \\
Fondo sin incendio & 0.0101 \\
\bottomrule
\end{tabular}
\captionof{table}{Comparación de energía de Haralick entre muestras reales del dataset.}
\end{center}

\begin{figure}[H]

{\centering \includegraphics[width=0.95\linewidth,height=\textheight,keepaspectratio]{figuras/sec08_glcm_comparacion_energia.png}

}

\caption{Comparación de energía GLCM entre las 3 muestras del dataset
propio.}

\end{figure}%

Las regiones de fuego presentan \textbf{mayor} energía que el fondo,
resultado contraintuitivo: la energía mide uniformidad de intensidades,
y las llamas generan grandes zonas contiguas de brillo saturado,
mientras que el fondo (con múltiples elementos de escena) dispersa más
las co-ocurrencias de intensidad.

\textbf{SIFT/ORB.} Al elevar \texttt{nfeatures} de 200 a 500, los puntos
clave detectados subieron de \textasciitilde142 a \textasciitilde165 por
imagen, y las coincidencias válidas de 107 a 122.

\subsection{Redes neuronales convolucionales y aplicaciones de IA
(Sesión
2)}\label{redes-neuronales-convolucionales-y-aplicaciones-de-ia-sesiuxf3n-2}

\textbf{Procesamiento en tiempo real.} El benchmark de tiempo por cuadro
(presupuesto de 33.33 ms para 30 FPS) mostró que tanto la configuración
a color (2.26 ms, 442.5 FPS) como en gris con resolución reducida (1.72
ms, 580.3 FPS) superan el umbral en este entorno, aunque la versión
optimizada mantiene una ventaja relevante para hardware con menor
capacidad de cómputo.

\textbf{CNN desde cero y profundidad adicional.} Se comparó la
arquitectura base (2 bloques Conv+Pooling) contra una versión con una
tercera capa \texttt{Conv2D(64,3)} (Tabla 2).

\begin{center}
\small
\begin{tabular}{@{}lcc@{}}
\toprule
\textbf{Métrica} & \textbf{Sin 3.\textsuperscript{a} capa} & \textbf{Con 3.\textsuperscript{a} capa} \\
\midrule
Accuracy & 83.00\% & 84.00\% \\
Falsos negativos & 22 & 13 \\
Falsos positivos & 12 & 19 \\
Recall (Fire) & 85.4\% & 91.4\% \\
Precisión (Fire) & 91.5\% & 87.9\% \\
\bottomrule
\end{tabular}
\captionof{table}{Comparación de arquitecturas CNN sobre el conjunto de prueba.}
\end{center}

\begin{figure}[H]

{\centering \includegraphics[width=0.95\linewidth,height=\textheight,keepaspectratio]{figuras/sec14_matriz_confusion_3capa.png}

}

\caption{Matriz de confusión de la CNN con tercera capa convolucional
(accuracy: 84.00\%).}

\end{figure}%

El accuracy global cambió poco (+1 punto), pero el tipo de error cambió
sustancialmente: los falsos negativos cayeron un 41\%, a costa de más
falsas alarmas --- una mejora relevante dado el mayor costo de no
detectar un incendio real.

\textbf{Regularización.} Dropout=0.5 combinado con \emph{data
augmentation} y \emph{early stopping} alcanzó 85.50\% de accuracy, con
curvas de validación más estables que el modelo sin regularizar.

\textbf{Transfer learning y fine-tuning.} La Tabla 3 resume el desempeño
de los cuatro modelos evaluados.

\begin{center}
\small
\begin{tabular}{@{}lcc@{}}
\toprule
\textbf{Modelo} & \textbf{Época parada} & \textbf{Accuracy} \\
\midrule
CNN desde cero & 7 (mejor: 2) & 84.00\% \\
Regularización combinada & 11 (mejor: 6) & 85.50\% \\
Transfer learning (MobileNetV2) & 11 (mejor: 8) & 91.00\% \\
Fine-tuning (últimas 10 capas) & 4 (mejor: 1) & 89.00\% \\
\bottomrule
\end{tabular}
\captionof{table}{Comparación final de las cuatro arquitecturas evaluadas.}
\end{center}

\begin{figure}[H]

{\centering \includegraphics[width=0.95\linewidth,height=\textheight,keepaspectratio]{figuras/sec20_comparacion_modelos.png}

}

\caption{Comparación final de accuracy entre las cuatro arquitecturas
evaluadas.}

\end{figure}%

El transfer learning con MobileNetV2 superó ampliamente a los modelos
entrenados desde cero, aprovechando el conocimiento preentrenado en
ImageNet. El fine-tuning de las últimas 10 capas no mejoró el resultado
del transfer learning puro en este dataset específico.

\textbf{Métricas de detección (IoU y mAP).} Un desajuste de +15 píxeles
(frente a +3) redujo el IoU de 0.8884 a 0.5656 (caída de 0.3228). Al
exigir un umbral de mAP más estricto (0.7 frente a 0.5, estándar tipo
COCO), la fracción de detecciones válidas cayó de 86.00\% a 60.00\% (−26
puntos).

\subsection{Hallazgos del proceso de
revisión}\label{hallazgos-del-proceso-de-revisiuxf3n}

La verificación cruzada entre resultados impresos y respuestas escritas
permitió corregir cinco inconsistencias: (1) valores de accuracy citados
como rangos genéricos en vez de las cifras reales obtenidas; (2)
adaptación necesaria del ejercicio de GLCM al no contar con las texturas
de referencia originales; (3) una interpretación que contradecía el
propio benchmark de FPS impreso; (4) ausencia del valor de referencia
sin la tercera capa convolucional, necesario para una comparación
cuantitativa válida; y (5) un valor de accuracy hardcodeado en el código
de comparación final que no reflejaba el resultado real de la CNN desde
cero.

\subsection{Conclusiones}\label{conclusiones}

Los resultados confirman la progresión metodológica esperada: los
descriptores clásicos ofrecen desempeño competitivo a bajo costo
computacional, pero el transfer learning logra el mejor resultado
global. El hallazgo más relevante desde una perspectiva de ingeniería es
que el accuracy global puede ocultar diferencias importantes en el tipo
de error cometido --- un matiz crítico en aplicaciones de seguridad como
la detección temprana de incendios.




\end{document}
